\Titular%
{divulgacion}%
{Kepler, as dúas primeiras leis e a revolución astronómica}%
{Iago Arsequell Rodríguez}%
{}%

\begin{multicols}{2}

\subsection*{Introdución}

A Revolución Científica foi un dos eventos chave da Historia. Este, por raro
que pareza, foi motivado por problemas non intrínsecos á Terra senón máis alá
desta. É que a pregunta de como é o que vai máis alá do ceo parécenos agora
trivial, que até un neno de primaria sabe máis ou menos a resposta. En cambio,
quixera amosar a complexidade e o esforzo requirido para solucionar esta
cuestión que por séculos foi un dos maiores misterios. Nesta tesitura, Kepler
foi quen de non só responder e nos dar un gran avance no coñecemento, senón un
xeito de enfrontarse ao dilema cunha mentalidade e métodos innovadores para a
época. Falarase de como atopou as súas dúas primeiras leis que desenvolveu no
seu libro \textit{Astronomia Nova}, libro de suma importancia na historia da
astronomía e da ciencia.

\subsection*{Proceso de dedución das dúas primeiras leis de Kepler}

Tycho Brahe fora un nobre danés e un dos maiores astrónomos antes da invención
do telescopio. Construíu un observatorio chamado Uraniborg, que se convertiu no
lugar cos instrumentos e medicións máis punteiras da época. Entre os seus
maiores logros áchanse medicións da posición dos planetas e ángulos destes
cunha precisión de dous arco minutos (un minuto de arco equivale ao ancho da
mina dun lapis vista a dous metros e medio de distancia). Foran medicións que
sen ter telescopios nin material máis avanzado resultan incribles e exténdense
vinte anos . Tras diferenzas co rei, abandona Dinamarca e acaba en Praga onde
sería astrónomo real do Sacro Imperio Romano Xermánico. Aquí é onde entra
Kepler. Brahe tiña os datos, mais non a capacidade para demostrar nada.
Persoalmente   tiña o seu modelo onde a Terra era o centro do universo e o Sol
orbita arredor da Terra coa diferenza de que o resto de planetas orbitan o Sol
(un modelo híbrido). Kepler, apaixonado pola idea dun universo xeométrico e
deseñado por Deus, quería demostrar que nos datos se obtiñan patróns
xeométricos e leis matemáticas que eran reflexo dunha vontade divina. En
resumo, un tiña os datos e o outro a capacidade de usalos.

Brahe reticente de darlle os datos a Kepler de golpe (e moitos outros
problemas) decidiu darlle os datos de Marte. Sería con este planeta co que
principalmente Kepler traballaría, o que é unha casualidade histórica. Marte é
un dos planetas con órbitas máis excéntricas do sistema solar polo que é máis
doado chegar  a resultados que co resto de planetas. Primeiro, repasemos uns
detalles da  astronomía kepleriana.

\begin{figure}[H]
    \centering
    \includegraphics[width=\linewidth]{revistas/005/imaxes/0a.png}
    \caption{Modelo ptolemaico.}
\end{figure}

A astronomía tradicional aristotélica pon á Terra no centro tal que os planetas
orbitan a velocidade constante. Este modelo non concorda en absoluto coas
observacións, como por exemplo co fenómeno do movemento retrógrado. Nas
observación antigas xa se observaba que os planetas se moven de oeste a leste
mais en certos tempos a velocidade diminúe (visto dende a Terra), paran e
inverten o seu movemento até desacelerar, parar e voltar ao movemento de antes
debuxando unha forma de lazo, este movemento non é igual sempre. Para este tipo
de fenómenos Ptolomeo desenrolou o seu modelo complexo que pretendía explicar e
predicir as observacións que se tiñan na súa época. Para entender o sistema,
describimos os trazos máis importantes.  Primeiro, a Terra é o centro do
universo. Os planetas móvense en círculos chamados epiciclos describindo un
círculo chamado deferente, no seu centro non se atopa a Terra, senón que esta
está desprazada unha distancia. O outro trazo máis importante é o ecuante que é
o punto oposto á Terra (isto lembra aos focos dunha elipse). Este punto
defínese como o punto dende o cal o movemento do centro do epiciclo é circular
uniforme. Isto explica tamén por que parece variar a velocidade dos planetas.
Todo este sistema que parece tan enrevesado se creou para poder predicir as
observacións astronómicas e as antigas concepcións (como o movemento circular
uniforme).

É posible que un modelo xeocéntrico describa a retrogradación dos planetas?
Ptolomeo ideou a súa teoría en boa parte para tentalo xustificar. Kepler, un
xeocéntrico convencido, demostrou que si, o fenómeno pódese explicar pola
diferenza de velocidade. Cando a Terra está afastada de Marte, véselle facendo
unha traxectoria recta. Pola contra, cando se aproxima a Terra a Marte (ao
estar máis cerca do Sol ten unha maior velocidade) a posición aparente de Marte
desacelérase producindo este fenómeno (na seguinte imaxe vese mellor).

\begin{figure}[H]
    \centering
    \includegraphics[width=\linewidth]{revistas/005/imaxes/0b.png}
    \caption{Retrogradación de Marte.}
\end{figure}

Unha dificultade engadida fora que a Terra non era un observador fixo, senón
que tamén se move. Sería moito máis sinxelo observar dende o Sol, e isto pódese
facer aproveitando as oposicións, momentos no que o Sol e o planeta están
aliñados nos lados opostos da Terra. Nese caso, por xeometría básica, o ángulo
entre o Sol e Marte é o mesmo visto dende a Terra que dende o Sol. Se se coñece
a distancia Terra-Sol, pódese usar perfectamente o Sol como referencia. Isto
permite contestar unha pregunta fundamental dende tempos antigos: é a
velocidade angular dos planetas uniforme? Pensadores tan sobresalintes coma
Platón ou Aristóteles estableceran que si. Kepler tomou coma referencia doce
oposicións, de 1580 a 1604. Se a velocidade angular fora uniforme e contando
que o ano marciano son 687 días:

\begin{equation}
    T_M=687\ \text{días} \Rightarrow \omega_M = \frac{360^\circ}{687 \text{días}}=0.524^\circ /\text{días}
\end{equation}

Cos datos observacionais, Kepler podía predicir onde debería de estar Marte na
próxima oposición. Calculando a posición das próximas o error acumulado era moi
superior aos dous arco minutos mencionados antes. É importante notar que esta
rigorosidade en ter errores aceptables é algo meritorio de Kepler nunha época
na que non se acostumaba. Isto permitiulle demostrar como o movemento angular
dos planetas non é constante.

\begin{figure}[H]
    \centering
    \includegraphics[width=\linewidth]{revistas/005/imaxes/1.png}
    \caption{Oposición Marte e ángulos do Sol e Tierra.}
\end{figure}

% \begin{figure}[H]
%     \centering
%     \includegraphics[width=\linewidth]{revistas/005/imaxes/2.png}
%     \caption{
%         Esquema coa formulación do problema da velocidade co error das
%         predicións.
%     }
% \end{figure}

O seguinte paso é construír un modelo de órbita. Se non é velocidade constante,
a velocidade variará na órbita. Isto derivou en recuperar a idea do ecuante.
Recuperalo provoca que, ao ter que conservar o movemento circular respecto do
ecuante, Marte móvase máis rápido cando está máis afastado del (e, por tanto,
está máis cerca do Sol) e máis lento cando está máis cerca (máis lonxe do Sol).
Kepler incorporaríao coidadosamente, construíndo o modelo do seguinte xeito.
Primeiro, Ptolomeo asume que a distancia do centro da órbita ao Sol ($e_{Sol}$)
e ao ecuante ($e_{ecuante}$) son iguais, Kepler asumiraos como variables. Logo,
debido a que a velocidade angular é constante, podemos tomar catro vectores do
ecuante a Marte cuns ángulos hipotéticos (atópanse dividindo o período orbital
entre os días que hai entre observacións) para comparalos cos catro vectores
dos ángulos do Zodíaco (vistos dende o Sol) que son os ángulos observacionais
(tomando só catro posicións).

\begin{figure}[H]
    \centering
    \includegraphics[width=\linewidth]{revistas/005/imaxes/3.jpg}
    \caption{Variables e resultados da \textit{Hipótesis Vicaria}.}
\end{figure}

Por un axuste de computación a forza bruta podemos achar o círculo único
(órbita superposta) onde intersecan os dous conxuntos de catro vectores. Este
método tan tedioso que levou moito tempo resultou  un éxito ao poñer o resto de
observacións xa que deu un error total de 2 arco minutos. Dito modelo chamado
\textit{Hipótesis Vicaria} (\textit{Vicaria} quere dicir \textit{sustitutiva})
pode semellar verdadeiro pero Kepler quería comprobalo antes.

Kepler desenrolou un método que con só simple trigonometría podía saber a
distancia do Sol a Marte. Os planos orbitais da Terra e Marte non son iguais
senón que están inclinados cun ángulo $\alpha\thickapprox1.85$º. Isto permite
facer o seguinte triángulo entre o Sol, Marte e a Terra a partir dunha
oposición onde a distancia vertical entre os planos orbitais sexa máxima.
Engadindo unha segunda oposición no lado oposto pódese facer un sistema de
dobre triángulo co que se pode, coñecendo a distancia Sol-Terra (coñecida por
Kepler e tendo un valor mínimo e máximo) e os ángulos do sistema, deducir a
distancia do Sol a Marte e ao centro da órbita.

\begin{figure}[H]
    \centering
    \includegraphics[width=\linewidth]{revistas/005/imaxes/4.png}
    \caption{Problema trigonométrico xerado polas dúas oposicións.}
\end{figure}

Kepler descubriu que comparando a distancia Marte-Sol preditas e utilizadas
polo modelo diverxían polo menos un 14\% e até un 42\% entre as calculadas coas
medicións directas. Isto foi un gran problema que fixo a Kepler dubidar cal
sería o erro, se talvez é o círculo perfecto que se supón como órbita ou o
ecuante. Se se modifica a \textit{Hipótesis Vicaria} para poñer as distancias
Marte-Sol calculadas, a distancia ao centro do ecuante e ao Sol é
aproximadamente igual. O mesmo que o modelo de Ptolomeo. A Kepler parecíalle
que funcionaba, mais na súa \textit{Hipótesis Vicaria} modificada cando se
calculaban as predicións agora devolvían un erro de oito arco minutos.

Estes oito minutos son clave para revolucionar a astronomía. Non é un gran
error. O propio modelo de Ptolomeo non baixaba dos dez arco minutos. Kepler non
estaba convencido cos erros xa que o das obervacións era na súa demasía de dous
minutos. No lugar de ignoralo, entendeu que necesitaba buscar novas causas e
sistemas e repensar todo o que fixera até agora, remeditar toda a astronomía.
Que fan que se movan os planetas? Ninguén pretendera seriamente entender isto.
Kepler imaxinouse ao Sol coma un gran imán que facía unha forza que chamou unha
<<\textit{especie inmaterial do corpo solar}>>. Esta forza sería inversa á
distancia como a luz que é máis débil [[[[como la luz que es más débil como más
dispersa es (es una duda si hubiera sabido la ley inversa cuadrática de la luz
si hubiera sido capaz de deducir la ley de gravitación).]]]] Cuanto máis
afastado da Terra, menor forza e menor velocidade. Este intento de explicación
física intúe unha posible relación matemática. Kepler sabía de como Arquímedes
demostrara a relación entre a circunferencia e o diámetro dividindo o círculo
entre triángulos cada vez máis pequenos. Un método que precedía ao uso de
infinitesimais  e que o mesmo Kepler pretendeu replicar. Se dividimos a órbita
en triángulos de base $d$ (a variación da posición do planeta) e a altura $r$
(a posición do Sol ao planeta), o área de cada triángulo é
$A=\frac{1}{2}bh=\frac{1}{2}dr$. A velocidade é inversamente proporcional á
posición $v=\frac{k}{r}=\frac{d}{t}$. Entón, o tempo é directamente
proporcional ao producto de $r$ e $d$ e inversamente a unha constante $k$ tal
que $t=\frac{rd}{k}$. Este método de triangulación é a segunda lei de Kepler
que postula que o planeta se move nun intervalo de tempo que barre áreas
iguais. Kepler non sabía como determinar a posición exacta en cada instante
porque a velocidade varía a cada instante algo que co cálculo moderno sería
sinxelo. Isto coñécese como o \textit{problema de Kepler} que motivaría o
desenvolvemento do cálculo posterior. Tras crear este método, Kepler
abandonaría o concepto de ecuante.

Se o aplicamos á \textit{Hipótesis Vicaria}, as predicións que nos devolven son
dun erro ínfimo en certas posicións e errores enormes noutras. Isto motivou a
Kepler a dar o golpe de graza á astronomía antiga: ás órbitas circulares. Se
modificaba a forma da órbita facendo que os lados  do círculo se acercaran máis
ao centro, os resultados melloraban. Kepler desevolveu un método moi orixinal
para saber a forma da órbita. Marte e a Terra teñen posicións diferentes porque
se están a mover. Se puidéramos fixarnos nunha posición, poderíamos facilmente
saber a posición do outro. Isto é precisamente o que conseguiu Kepler. Se
tomamos unha posición de Marte, a Terra estará nunha posición determinada. Nun
período marciano, Marte volverá á mesma posición, pero a Terra estará nunha
distinta. Se recopilamos suficientes posicións da Terra, podemos saber a forma
da Terra de maneira aproximada xa que a sucesión de posicións revelan a órbita.

\begin{figure}[H]
    \centering
    \includegraphics[width=\linewidth]{revistas/005/imaxes/5.png}
    \caption{
        Debuxo do propio Kepler onde mostra a Marte na mesma posición e a
        posición da Terra despois de varios períodos de Marte.
    }
\end{figure}

A conclusión á que chegou Kepler foi que a forma da órbita debe ser algunha
clase de forma oval. Mais, que tipo de óvalo? El pensaba que unha forma
semellante á dun ovo, no entanto, hai multitude de óvalos entre eles a elipse
que Kepler desestimou por crer que era demasiado simple e que en tal caso
«tería sido»  moi evidente para os gregos. Kepler desenrolou un método de
epiciclos (tanto tempo quitándoos para rematar usándoos resulta irónico) para
simular a órbita que lle resultaba na forma dun ovo pero que non servía para
predicir as observacións. Comentáramos antes que daba mellores resultados se
ambos lados opostos do círculo se encollían cara o interior, comparados á
curvatura resultado da triangulación. A máxima distancia entre o círculo e esta
curvatura é de 0’0049. Kepler desenrolara paralelamente unha ecuación que lle
daba o ángulo de Marte respecto ao vector posición entre o centro e o Sol que
Kepler chamaría \textit{ecuación óptica}. Por pura casualidade, tras analizar
esta ecuación deuse de conta que en 5’30º o seu valor máximo, cando o ángulo do
centro e do Sol é de 90º, se tomamos a secante de 5’30º, o resultado daba
1’0049. Xusto a distancia máxima máis un! Se modifica o seu modelo e axusta o
seu método de epiciclos para facer que a distancia Marte-Sol foran
proporcionais á secante da ecuación óptica, o seu modelo encaixaría coas
observacións. Este complexo modelo xeométrico trazaba unha elipse perfecta.
Aínda así, Kepler equivocouse inicialmente e frustrado desestimou este modelo.
Emocionado por atopar novas ideas, pensou sobre figuras xeométricas. Ao probar
a elipse deuse de conta que funcionaba e que xa a descobrera antes co devandito
modelo. Tras catro largos anos, descubriuse a primeira lei de Kepler (nótese
que a segunda foi anterior á primeira).

\begin{figure}[H]
    \centering
    \includegraphics[width=\linewidth]{revistas/005/imaxes/6.png}
    \caption{
        Esquema do problema formulado coa curvatura de círculo perfecto e
        observada e a ecuación óptiy el secante .
    }
\end{figure}

\subsection*{Conclusión.}

A principal motivación deste artigo é a de amosar a relevancia de todo este
desenvolvemento. Na época de Kepler, os métodos científicos eran practicamente
inexistentes. El foi quen de deixar dun lado a tradición e preconcepcións en
pos de desenvolver unha cadea de razóns lóxicas que xustificaban as dúas leis
de Kepler. É moi difícil esaxerar como de significativas foron estas leis e a
futura lei de gravitación (motivada principalmente polos resultados de Kepler)
para a historia e a ciencia. Comparemos a absoluta confusión e incompletitude
do coñecemento humano de épocas anteriores coa claridade e simplicidade destas
leis (sobre todo a da gravitación) que establecen de xeito sinxelo como se
moven os astros, algo de aparente gran complexidade e misticismo. Esta é a
razón do triunfo da ciencia a futuro, a esperanza de que todos os fenómenos do
mundo podían ter leis tan belas e simples.

% \nocite{historia_marte}
% \nocite{bizarre_mars}
\nocite{kepler_disc}
\nocite{shoulders}
\nocite{koestler_2017}
\nocite{kepler_2015}

\printbibliography

\end{multicols}
